\section{Discussion}
\label{sec:discussion}

Our results establish that heavy-tailed self-regularization theory extends to Vision Transformer attention mechanisms, while revealing that training provenance---not architecture---dominates cross-model variance.

\subsection{Key Findings}

\textbf{Extension of HT-SR to Attention Mechanisms.} We successfully computed heavy-tailed exponents for 53 ViT models using the Hill estimator on attention weight matrices. All processed models exhibited heavy-tailed distributions with $\alpha$ values in the [2.20, 18.92] range.

\textbf{Training Origin Dominates Variance.} The 12$\times$ variance reduction through family stratification reveals that training provenance is the primary variance driver. Models sharing training origin show consistent $\alpha$ values regardless of architectural variants.

\textbf{Task-Specific Fine-Tuning Creates Outliers.} Models fine-tuned for specialized tasks exhibit extreme $\alpha$ values ($>$10), reflecting genuine distributional shifts.

\subsection{Limitations}

\textbf{Sample Size Below Target.} We processed 53 of 100 target models (53\% success rate). The pattern is statistically robust, but additional models would strengthen confidence.

\textbf{Single Family Tested.} Within-family variance reported only for google/vit-*. Other families had insufficient samples for reliable estimation.

\textbf{Downstream Hypotheses Blocked.} H-M1, H-M2, H-M3 testing cross-architecture prediction transfer remain untested pending hypothesis revision.

\subsection{Broader Impact}

This work advances understanding of weight-space analysis for model selection. Weight-based metrics could be gamed if publishers manipulate distributions without improving quality. We recommend combining weight metrics with inference-based validation for high-stakes applications.
